\documentclass[11pt]{article}
\usepackage[margin=1in]{geometry}
\usepackage[T1]{fontenc}
\usepackage{amsmath,amssymb,mathtools,bbm}
\usepackage{booktabs,longtable,array,enumitem}
\usepackage{listings,xcolor,hyperref}
\sloppy
\hypersetup{colorlinks=true, linkcolor=blue!50!black, urlcolor=blue!50!black}
\definecolor{cbg}{rgb}{0.97,0.97,0.97}
\lstset{basicstyle=\ttfamily\scriptsize, backgroundcolor=\color{cbg}, showstringspaces=false,
  breaklines=true, frame=single, framerule=0pt, columns=fullflexible}
\newcommand{\code}[1]{\texttt{#1}}
\newcommand{\file}[1]{\texttt{#1}}
\setlist{nosep, leftmargin=1.4em}

\title{AutoBayes brms-parity loop: Tier A review}
\author{Status at the Tier A boundary, prepared for Andrew}
\date{7 October 2026}

\begin{document}
\maketitle

\noindent This is the tier-boundary review the approved plan calls for: what landed, what parked, where the library
knowingly differs from brms, how each claim was checked, and what needs your decision before Tier B starts.
Section~\ref{sec:decisions} is the part that needs you; the rest is the evidence behind it.

\section{What ``merged'' means in this loop}
\label{sec:meaning}

Every feature is accepted by comparison with brms itself, never by sampling. For a model brms accepts, the oracle
(\file{oracle/R/gen\_fixture.R}) records, once, deterministically: the Stan data (\code{standata()}), the log density
and its gradient at eight seeded points of the unconstrained space (\code{rstan::log\_prob}, with and without the
Jacobian), and \code{posterior\_linpred}/\code{posterior\_epred} at those points. That record is a frozen fixture. A
feature merges when the AutoBayes model built from the same data reproduces it:

\begin{itemize}
\item G1 design: design matrices, response, trials, offsets and group coding equal \code{standata} to $10^{-12}$;
\item G2 density: log density, its prior part and every gradient coordinate equal brms to $10^{-8}$ / $10^{-7}$
  relative (relative to $\max(1, |a|, |b|)$, so at a log density near $-120$ the floor is about $10^{-6}$ absolute);
\item G3 derived quantities and G4 predictions to $10^{-9}$; G5 the JAX density equals the PyTensor one;
  G6 \code{ab.default\_priors()} reproduces brms's \code{default\_prior()} table string for string (strict since A07);
\item G7 the whole non-sampling suite locally, and the full fast suite in hosted CI on the pull request.
\end{itemize}

\noindent What this does \emph{not} establish: that NUTS converges on these targets, or that intervals are calibrated.
It establishes that the Python model \emph{is} brms's target, coordinate by coordinate. Nothing in the gates samples:
the only sampling in Tier A was the existing short sampling tests that hosted CI runs on every pull request, and the
advisory 120-second smoke (G8) was never invoked.

\section{What landed}
\label{sec:landed}


\small
\begin{center}\begin{tabular}{@{}>{\raggedright\arraybackslash}p{0.16\textwidth}p{0.05\textwidth}>{\raggedright\arraybackslash}p{0.35\textwidth}>{\raggedright\arraybackslash}p{0.25\textwidth}>{\raggedright\arraybackslash}p{0.08\textwidth}@{}}
\toprule
Feature & PR & What a user can now write & Frozen brms cases it passes & Diverg. \\
\midrule
A01 independent slopes & \#29 & \code{(1 + x [+ w] || g)}, any numeric slopes; \code{sd=\{"g": \{"Intercept": ..., "x": ...\}\}} (your API decision) & \code{gauss\_\allowbreak{}slope\_\allowbreak{}independent}, \code{\_\allowbreak{}two\_\allowbreak{}groups}, \code{\_\allowbreak{}two\_\allowbreak{}slopes}, \code{bernoulli\_\allowbreak{}slope\_\allowbreak{}categorical}, \code{poisson\_\allowbreak{}slope\_\allowbreak{}independent} & D8 \\
A02 correlated slopes & \#35 & \code{(1 + x | g)} with \code{cor=\{"g": ab.LKJ(eta)\}}; $L$ sampled in Stan's own Cholesky-correlation coordinates & \code{gauss\_\allowbreak{}slope\_\allowbreak{}correlated} & D20, D21 \\
A03 grouping syntax & \#31 & \code{(1 | g:h)}, \code{(1 | g/h)}, brms's \code{g\_\allowbreak{}h} labels & \code{gauss\_\allowbreak{}nested\_\allowbreak{}colon}, \code{\_\allowbreak{}slash} & D11, D12 \\
A04 interactions & \#34 & \code{x:w}, \code{z:h} (two-way, main effects present) & \code{gauss\_\allowbreak{}interaction\_\allowbreak{}numeric\_\allowbreak{}numeric}, \code{\_\allowbreak{}categorical\_\allowbreak{}categorical} & D14, D15 \\
A05 transforms & \#34 & \code{log(x)}, \code{exp(x)}, \code{sqrt(x)}, \code{I(x\^{}k)}, \code{I(x * k)}, \code{I(x + k)}, \code{I(x - k)} & \code{gauss\_\allowbreak{}transforms} & D16 \\
A06 offset & \#31 & \code{offset(col)}, \code{offset(log(col))} & \code{poisson\_\allowbreak{}offset} & D10 \\
A08 negative binomial & \#30 & \code{ab.negbinomial}, \code{Priors(shape=...)} & \code{negbinomial\_\allowbreak{}varint} & -- \\
A09 Gamma & \#30 & \code{ab.gamma}, log link & \code{gamma\_\allowbreak{}varint} & D9 \\
A10 lognormal & \#32 & \code{ab.lognormal}; \code{kind="mean"} is brms's $e^{\mu+\sigma^2/2}$ & \code{lognormal\_\allowbreak{}varint} & -- \\
A11 Beta & \#32 & \code{ab.beta}, \code{Priors(phi=...)} & \code{beta\_\allowbreak{}varint} & D13 \\
A12 links & \#36 & \code{link="probit"} for Bernoulli and binomial; cloglog parked (Section~\ref{sec:park}) & \code{bernoulli\_\allowbreak{}probit}, \code{binomial\_\allowbreak{}probit} & D17--D19 \\
A13 learned $\nu$ & \#36 & \code{ab.student(..., nu=None)} with \code{Priors(nu=...)}, truncated at 1 & \code{student\_\allowbreak{}learned\_\allowbreak{}nu} & D22 \\
A07 default priors & \#37 & \code{ab.default\_priors(formula, data, family=...)}: brms's defaults on request, data-dependent intercept location and scales reproduced exactly & G6 on all 38 ungated cases; 5 new cases frozen without priors & D1, D2, D23 \\
\bottomrule
\end{tabular}\end{center}
\normalsize



\noindent Infrastructure behind these (Phase 0, all merged): \#22 an automatic \code{sampling} marker so the
deterministic suite runs alone; \#23 the oracle and its first 14 fixtures; \#24 the loop's rules, ledger and gate
tools; \#25 a family registry (no \code{family ==} branches remain; a lint test forbids new ones); \#27 an explicit
term model (\code{PopTerm}, \code{GroupTerm}, per-coefficient group blocks), which is what let slopes, grouping
syntax and interactions land as small additions rather than rewrites. Fixtures for every feature were frozen before
any implementation existed (\#26, \#28, \#33) and ran as strict expected failures until their branch removed the gate.

\paragraph{How close.} The gates only say ``within tolerance''. The table below says by how much: for every
ungated case, the worst relative error against brms over the eight points (log density with and without the Jacobian;
every gradient coordinate; linear predictor and expected value on fitted and new rows, with and without group effects).
Most sit at $10^{-16}$--$10^{-15}$, i.e.\ the same number up to rounding. The exceptions follow the special
functions involved: the cases with Student-$t$-type densities (the Student-$t$ likelihood; half-Student-$t$ and half-Cauchy priors,
which includes every case run under brms's default priors) have log-density differences of $2\times10^{-11}$ to
$2.3\times10^{-10}$, and the families whose
gradients involve digamma (Beta, Gamma, negative-binomial shape, learned $\nu$) have gradient differences of
$4\times10^{-11}$ to $9.4\times10^{-10}$; for the last, the A13 work traced it to PyTensor's C digamma (about
$2\times10^{-10}$ accurate, where JAX's agrees with brms to $3\times10^{-15}$). I have not traced the Student-$t$ density
difference further. Even the worst cases are 40, 100 and $10^{6}$ times inside the $10^{-8}$, $10^{-7}$ and $10^{-9}$
tolerances. (Measured on the A07 branch, which contains everything in Tier A; 38 cases.)

\scriptsize\setlength{\tabcolsep}{4pt}
\begin{center}\begin{tabular}{@{}lllrrr@{}}
\toprule
case & family & coords & log density & gradient & predictions \\
\midrule
\code{bernoulli\_interaction\_two\_groups} & bernoulli & 18 & 3.0e-16 & 8.9e-16 & 6.3e-16 \\
\code{bernoulli\_probit} & bernoulli & 9 & 3.3e-16 & 4.5e-15 & 3.6e-16 \\
\code{bernoulli\_slope\_categorical} & bernoulli & 18 & 2.8e-16 & 1.0e-15 & 4.4e-16 \\
\code{bernoulli\_slope\_default\_priors} & bernoulli & 17 & 1.0e-10 & 6.7e-16 & 4.0e-16 \\
\code{bernoulli\_varint} & bernoulli & 9 & 3.2e-16 & 1.2e-15 & 2.2e-16 \\
\code{beta\_varint} & beta & 10 & 9.6e-16 & 2.7e-10 & 4.1e-16 \\
\code{binomial\_poponly\_categorical} & binomial & 4 & 6.7e-16 & 5.0e-16 & 4.6e-16 \\
\code{binomial\_probit} & binomial & 9 & 9.4e-15 & 2.3e-13 & 1.0e-15 \\
\code{binomial\_varint} & binomial & 9 & 6.5e-16 & 1.3e-15 & 4.4e-16 \\
\code{gamma\_varint} & gamma & 10 & 7.3e-16 & 3.0e-10 & 4.6e-16 \\
\code{gauss\_categorical} & gaussian & 12 & 4.6e-11 & 1.4e-15 & 4.4e-16 \\
\code{gauss\_interaction\_categorical\_categorical} & gaussian & 13 & 2.7e-16 & 1.3e-15 & 3.3e-16 \\
\code{gauss\_interaction\_numeric\_numeric} & gaussian & 12 & 4.7e-16 & 3.7e-15 & 6.4e-16 \\
\code{gauss\_interaction\_two\_groups} & gaussian & 19 & 3.6e-16 & 4.4e-15 & 3.3e-16 \\
\code{gauss\_intercept\_only} & gaussian & 9 & 4.0e-16 & 7.3e-16 & 0.0e+00 \\
\code{gauss\_nested\_colon} & gaussian & 30 & 2.9e-16 & 1.7e-15 & 4.4e-16 \\
\code{gauss\_nested\_slash} & gaussian & 31 & 2.7e-16 & 1.4e-15 & 4.4e-16 \\
\code{gauss\_no\_intercept} & gaussian & 5 & 5.2e-16 & 4.0e-15 & 0.0e+00 \\
\code{gauss\_poponly} & gaussian & 4 & 4.3e-16 & 7.7e-15 & 3.7e-16 \\
\code{gauss\_slope\_correlated} & gaussian & 18 & 3.8e-16 & 3.2e-15 & 5.6e-16 \\
\code{gauss\_slope\_independent} & gaussian & 17 & 2.1e-16 & 1.3e-15 & 4.4e-16 \\
\code{gauss\_slope\_two\_groups} & gaussian & 23 & 2.1e-16 & 1.6e-14 & 8.9e-16 \\
\code{gauss\_slope\_two\_slopes} & gaussian & 25 & 4.5e-16 & 7.2e-15 & 7.8e-16 \\
\code{gauss\_transforms} & gaussian & 11 & 5.9e-16 & 4.6e-15 & 3.9e-16 \\
\code{gauss\_varint} & gaussian & 10 & 3.5e-16 & 9.1e-16 & 2.3e-16 \\
\code{gauss\_varint\_default\_priors} & gaussian & 10 & 3.2e-11 & 1.5e-14 & 2.3e-16 \\
\code{gauss\_varint\_uncentered} & gaussian & 10 & 3.3e-16 & 3.4e-15 & 2.2e-16 \\
\code{lognormal\_varint} & lognormal & 10 & 3.7e-16 & 7.4e-16 & 9.7e-16 \\
\code{negbinomial\_default\_priors} & negbinomial & 10 & 2.0e-11 & 1.6e-10 & 4.4e-16 \\
\code{negbinomial\_varint} & negbinomial & 10 & 6.7e-16 & 3.8e-11 & 4.9e-16 \\
\code{poisson\_categorical\_default\_priors} & poisson & 11 & 2.1e-11 & 1.0e-15 & 5.6e-16 \\
\code{poisson\_offset} & poisson & 9 & 4.9e-16 & 2.6e-15 & 8.8e-16 \\
\code{poisson\_poponly\_categorical} & poisson & 4 & 3.0e-16 & 1.0e-15 & 4.3e-16 \\
\code{poisson\_slope\_independent} & poisson & 16 & 2.7e-16 & 2.3e-15 & 7.9e-16 \\
\code{poisson\_varint} & poisson & 9 & 6.1e-16 & 1.1e-15 & 4.3e-16 \\
\code{student\_default\_priors} & student & 11 & 2.2e-11 & 1.5e-10 & 3.8e-16 \\
\code{student\_learned\_nu} & student & 11 & 3.8e-16 & 9.4e-10 & 3.1e-16 \\
\code{student\_varint} & student & 10 & 2.3e-10 & 2.5e-15 & 2.2e-16 \\
\bottomrule
\end{tabular}\end{center}
\normalsize


\section{The one park: brms's cloglog link}
\label{sec:park}

The probit link landed for Bernoulli and binomial models. The complementary log-log link did not, and the reason is a
defect in brms's arithmetic rather than in ours, so it is worth reading through once.

With $\theta = 1 - \exp(-e^{\eta})$, a row with $y = 0$ contributes $\log(1-\theta) = -e^{\eta}$ to the log density.
Its derivative in $\eta$ is exactly $-e^{\eta}$. Stan instead evaluates \code{bernoulli\_lpmf(0 | inv\_cloglog(eta))},
which forms $1-\theta$ \emph{by subtraction} from a stored $\theta$.

\paragraph{The numbers at the frozen point.} At the sixth of the eight seeded points of \file{bernoulli\_cloglog}
(index~5 in the fixture), one $y = 0$ row has $\eta = 3.008$. Then $e^{\eta} = 20.25$ and
$1-\theta = e^{-20.25} \approx 1.6\times 10^{-9}$, while $\theta \approx 1 - 1.6\times10^{-9}$ is stored to about
$1.1\times 10^{-16}$ absolute. Subtracting leaves $1-\theta$ with up to $1.1\times10^{-16}/1.6\times10^{-9} \approx
7\times 10^{-8}$ relative error, and that row's gradient term is multiplied by $e^{\eta} \approx 20$. The fixture
records the coefficient gradient as $-1.0055917217903116$; the exact value (50-digit arithmetic from the fixture's own
data and parameters) is $-1.0055921064960804$, a relative error of $3.8\times10^{-7}$, above the $10^{-7}$ tolerance.
A float64 imitation of Stan's subtraction reproduces brms's number to $7\times10^{-15}$, and regenerating the fixture
from R reproduces it exactly, so this is deterministic brms behaviour, not a fixture fault. AutoBayes's stable form
($\log(1-\theta) = -e^\eta$ directly) is within $4.6\times10^{-15}$ of exact at that point; at the other seven points
the two agree to $2.7\times10^{-13}$.

\paragraph{Options.} (a) Keep it parked: \code{link="cloglog"} is refused with this reason, the stable form stays in
\file{likelihoods.py} behind tests (recommended). (b) Reproduce brms's subtraction so the gate passes: that copies a
loss of accuracy into the library on purpose. (c) Accept cloglog against a closed-form reference instead of brms for
this one link: possible, but it is the first feature that would not be oracle-accepted.

\section{Where AutoBayes knowingly differs from brms}
\label{sec:divergences}

Every difference is logged in \file{docs/brms-parity/divergences.md}; fixtures were never edited to hide one.
Most are deliberate restrictions of scope with a clear refusal message; a few are numerically better forms of the same
density that agree with brms to within the tolerances on every frozen point.

\small
\begin{center}\begin{tabular}{@{}p{0.05\textwidth}>{\raggedright\arraybackslash}p{0.14\textwidth}>{\raggedright\arraybackslash}p{0.28\textwidth}>{\raggedright\arraybackslash}p{0.28\textwidth}>{\raggedright\arraybackslash}p{0.13\textwidth}@{}}
\toprule
& area & brms & AutoBayes & status \\
\midrule
D1 & priors & default priors for every class & every prior explicit; a class the model lacks is refused & kept; brms's defaults on request (A07) \\
D2 & priors & intercept-only model needs no \code{b} prior & still asks for a vacuous \code{b} prior & resolved (A07): no \code{b} prior needed \\
D3 & factor levels & R factor order (locale-dependent) & levels sorted as strings & handled: matched by content \\
D4 & binomial & keeps zero-trial rows & drops and reports them & kept \\
D5 & binomial means & \code{posterior\_epred} is $n p$ & \code{kind="mean"} gives $p$ and \code{expected} $= np$ & documented \\
D6 & new levels & resamples by default & refuses by default; two explicit policies & documented \\
D7 & Bernoulli & \code{bernoulli\_logit\_lpmf} & gradient-stable softplus form & equal on fixtures \\
D8 & slopes & slope without its population term allowed & slope must also be a population term & kept \\
D9 & Gamma & inverse link by default & log link only & kept \\
D10 & offset & any expression, several offsets & \code{offset(col)}, \code{offset(log(col))}, one & kept \\
D11 & grouping & three or more columns combined & two columns & kept \\
D12 & grouping & duplicate partitions; colliding \code{\_} labels merge & both refused & kept \\
D13 & Beta & $1-\mu$ by subtraction (0 beyond $\eta\approx37$) & $1-\mu$ as $\mathrm{expit}(-\eta)$ & equal on fixture \\
D14 & interactions & missing main effects, any order & two-way, both main effects & kept \\
D15 & interactions & offset or transformed operands & main-effect columns only & kept \\
D16 & transforms & any R expression & seven forms of one column & kept \\
D17 & probit & $\Phi$ saturates beyond $-37.5$ / $8.25$ & stable $\log\Phi$ & equal on fixtures \\
D18 & cloglog & $1-\theta$ by subtraction & refused with the reason & parked \\
D19 & other links & cauchit; links of other families & logit and probit for Bernoulli/binomial & kept \\
D20 & correlations & correlation IDs, \code{gr()} & one matrix per term; IDs and \code{gr()} refused & kept \\
D21 & correlations & a scalar per coefficient pair & whole matrix \code{cor|g} and factor \code{L|g} & names only \\
D22 & Student-$t$ $\nu$ & any prior, renormalised; default \code{gamma(2, 0.1)} & Gamma, Exponential, half-normal, half-Cauchy & kept \\
D23 & default priors & a fixed Student-$t$ $\nu$ still listed with \code{gamma(2, 0.1)} & no $\nu$ prior when $\nu$ is fixed & documented \\
\bottomrule
\end{tabular}\end{center}
\normalsize

\noindent ``Kept'' means a deliberate restriction with a refusal that names the construct and says what to write
instead; ``equal on fixtures'' means a numerically safer form of the same density, within the tolerances on every
frozen point.


\section{A07: brms's default priors}
\label{sec:a07}

Explicit priors remain the default: nothing is filled in silently (D1). A07 adds brms's defaults \emph{on request}:
\code{ab.default\_priors(formula, data, family=...)} returns the \code{Priors} brms would use, each marked as a default
in \code{describe()} and in fit records. The hard part is that several of brms's defaults depend on the data. For the
intercept brms uses $\text{student\_t}(3, m, s)$ with $m$ the median of the response on the link scale and
$s = \max(2.5, \text{MAD})$, both rounded and printed by R's own rules; for example \code{poisson\_offset}'s table reads
\code{student\_t(3, -0.978004218445523, 2.5)}. These are transcribed from brms and R and reproduced exactly, so gate G6
now compares brms's \code{default\_prior()} table with the library's row by row and string by string, and passes on
all 38 ungated fixtures; a one-unit change in the 15th digit fails it (a harness test). Five new fixtures were frozen
with no priors at all on the brms side, so G2 also checks the densities \emph{under} the defaults: a flat \code{b},
the data-dependent intercept, the bounded half-Student-$t$ scales, \code{inv\_gamma(0.4, 0.3)} for the negative
binomial shape, \code{lkj(1)}, and \code{gamma(2, 0.1)} above 1 for a learned $\nu$.

Beyond the fixtures (scratch checks, not committed), the transcription matched R on about 340{,}000 values of
\code{round}, \code{as.character}, \code{median}, \code{mad}, \code{mean} and \code{log}, and a differential test
against \code{brms::default\_prior()} matched 329 of 330 synthetic configurations (the other has a constant Bernoulli
response, which the constructors refuse). That test is what found brms's rule for naming an interaction
(\code{x:z + z + x} names the term \code{zq:x}). Side effect: an intercept-only formula no longer needs a vacuous
\code{b} prior (D2 resolved). Found in passing and not yet fixed: the order of two lines in a learned-$\nu$ model's
\code{describe()} depends on Python's hash seed (PyMC orders those two value variables by hash); the target identity
is unaffected.


\section{How the claims were checked}
\label{sec:verification}

\begin{itemize}
\item Each feature: the gate tool on the exact commit (immutability of fixtures and tolerances, the oracle tests, the
  full non-sampling suite), recorded with the commit SHA in the ledger; then both hosted CI runs on the pull request;
  then squash-merge as andrew-dou. Changes to \file{tests/oracle/adapter.py} and \file{gates.py} were read in review
  for weakening every time; all were extensions (new parameter kinds), and the correlated-slope change replaced a
  blanket refusal with an agreement check that is itself tested both ways.
\item Each refactor (registry, term model) and each feature: every previously ungated oracle-case model's target
  identity, \code{describe()} text, names and coefficients were snapshotted before the change and compared after; all
  identical in every case.
\item Reproducibility at the boundary: \code{oracle/regen.py check --all} regenerated the 43 frozen fixtures on \code{main} from R
  and compared them with the committed ones; the six frozen since were checked the same way: all 49 identical.
\end{itemize}

\paragraph{Limits I know of.} Correlated slopes are checked against brms for two coefficients only (the one fixture);
three coefficients are checked against closed forms, and Stan's LKJ normalising constant for three or more
coefficients with $\eta \ne 1$ is not independently confirmed. The special-function differences above (at most $9.4\times10^{-10}$) are inside
tolerance today but would be the first to matter if a tolerance were ever tightened.

\section{Process and cost}
\label{sec:cost}

Each pair of features was implemented by one agent working in its own git worktree from a written brief, then
gated, reviewed and merged by me (squash, as andrew-dou, no AI attribution). Token counts reported by the ten agents
that completed: P3 0.31M, P4 0.33M, A01 0.45M, A08+A09 0.39M, A06+A03 0.50M, A10+A11 0.46M, A04+A05 0.62M, A02 0.36M,
A12+A13 0.58M, A07 0.59M, about 4.6M in total. Not included: four agents cut off by limits before reporting a count,
and my own orchestration. Each agent ran 34--83 minutes; the non-sampling suite (about 1{,}300 tests) takes 6--17
minutes on this machine depending on load.

Two interruptions. On 6 October the included usage ran out; state was persisted and the work resumed after the reset.
On 7 October the Fable 5.1 monthly spend limit stopped two agents mid-work; their uncommitted work was backed up,
carried onto the current \code{main}, and finished after you switched the session to Opus 5.5, with agents pinned to
Opus.


\section{Decisions before Tier B}
\label{sec:decisions}

\begin{enumerate}
\item \textbf{The cloglog park} (Section~\ref{sec:park}). I recommend keeping it parked: option (a). Options (b) and (c)
  are both available if you want the link.
\item \textbf{API changes already merged that you may want to look at.} None breaks an existing call, but they change
  what some calls mean:
  \begin{itemize}
  \item \code{ab.student}'s \code{nu} now defaults to \code{None}, meaning ``learn it'', which requires
    \code{Priors(nu=...)}. A call without \code{nu} used to fail as a missing argument; it now fails with a message
    saying what to pass.
  \item \code{(1 + x | g)} now means correlated effects (brms's meaning); it used to be refused. \code{(1 || g)} is
    accepted as \code{(1 | g)}.
  \item \code{Priors} gained \code{shape}, \code{phi}, \code{cor} and \code{nu}, all after the four original fields,
    so positional calls keep their meaning. Bernoulli and binomial models gained \code{link=}.
  \end{itemize}
\item \textbf{Tier B.} Ten brms fixtures are already frozen for its six features: distributional \code{sigma ~ z},
  \code{cens()}/\code{trunc()}/\code{weights()}, \code{mo()}, \code{s()}, cumulative and categorical families,
  zero-inflated and hurdle Poisson. At Tier A's rate that is roughly 1.5--2M agent tokens over three two-feature
  agents. I expect \code{s()} to park: brms builds the spline basis with mgcv, which nothing in the pinned stack
  reproduces, so it would mean porting mgcv's thin-plate basis (my estimate; not yet attempted). I would start with
  the addition terms and the zero-inflated/hurdle families, which reuse the most existing machinery. Go ahead?
\item \textbf{Cleanup.} Sixteen local worktrees under \file{AutoBayes-worktrees/} belong to merged or closed
  branches: \file{parity-p1}, \file{-p2}, \file{-p3}, \file{-p4}, \file{-a01}, \file{-a01-fixtures},
  \file{-tierA-fixtures}, \file{-tierB-fixtures}, \file{-fam1}, \file{-fam2}, \file{-fam3}, \file{-form1},
  \file{-form2}, \file{-a02}, \file{-a07}, and \file{m3b1-varying-slope} (the closed PR \#8). Their branches stay on
  GitHub.
  \file{group-exclusion} holds the virtual environment every worktree uses and stays; \file{parity-loopstate} is the
  merge-queue checkout of \code{main}; your original checkout \file{projects/AutoBayes} is untouched. I will remove
  the sixteen only on your yes.
\end{enumerate}

\noindent Also in this pass: a documentation pull request that adds a current-status section to the README and
SKILL.md (they still describe slopes, \code{g/h}, interactions and transforms as refused) without removing the older
milestone text, and corrects two stale divergence references.


\end{document}
